\subsubsection{Linear Differential Operator}

\paragraph{The Formal Linear Differential Operator}

We define the differentian opeartor as:

\begin{mystthmdefinition}{Differentiation Opeartor}{definition-lineardifferentialoperator-1}Let $D:L^2[a,b]\to L^2[a,b]$ be the mapping defined by:
\begin{equation}
D(u)=\frac{d}{dx}u
\end{equation}
for all $u\in L^2[a,b]$

\end{mystthmdefinition}
Now we can define the \textit{formal linear operator} by some opeartor algebra:

\begin{mystthmdefinition}{Formal linear differential operator}{definition-lineardifferentialoperator-2}Let $p(x)$ be any polynomial over $\mathbb{C}$ of positive degree:
\begin{equation}
p(x)=a_\nu x^\nu
\end{equation}
where $a_{\nu}$ are functions of $x$. Define $L:=p(D)$ as the \textbf{formal linear differential operator} that correspond to the \textbf{auxillary polynomial $p$}:
\begin{equation}
L=a_\nu D^{\nu}
\end{equation}
The \textbf{order} of $L$ is given by the order of $p$.

\end{mystthmdefinition}
We call it \textit{formal} since we want to contrast with \textit{concrete} linear differential opeartor which is a formal opeartor endorsed with boundary conditions(This convention is used in \textit{Mathematical Physics by Michael Stone and Paul Goldbart}). We can think of the difference between them as:

\begin{itemize}
\item The \textbf{nullspace of formal opeartor} is the space of general solution of homogenous diffferential equation
\item The \textbf{nullspace of concrete operator} is the space of particular solution to homogenous differential equation with boundary conditions
\end{itemize}

We'll explore the nullspace of formal operators with a bit more detail now.

\subparagraph{Nullspace of formal operator}

We first observe something very nice about solutions to homogenous differential equations: they belong to $C^{\infty}$(infinitely differentiable).

\subparagraph{The Bootstrapping Lemma}

\begin{mystthmlemma}{}{lemma-lineardifferentialoperator-3}Let $L$ be a formal operator of order $n$. Suppose the coefficients of the auxillary polynomial are $k$ times continously differentiable($a_{\nu}\in C^k$) then $\forall y\in \mathrm{Null}(L)$, $y$ is at least $k+n$ times contiously differentiable.

\end{mystthmlemma}
\begin{mystthmproof}{}{proof-lineardifferentialoperator-4}Suppose $y\in \mathrm{Null}(L)$, note trivially $y\in C^n$. Then:
\begin{equation}
\begin{aligned}
a_0y+a_1y^{(1)}+a_2y^{(2)}...+a_n y^{(n)} &=0 \\
y^{(n)} &= -(b_0y+b_1y^{(1)}+b_2y^{(2)}...\underbrace{b_{n-1}{y}^{n-1}}_{\in C^{\min(k, 1)}})
\end{aligned}
\end{equation}
Where $b_{\nu}=a_{\nu}/a_{n}$. Now $b_{\nu}\in C^k$ by closedness under polynomial division. Also $y^{(\nu)} \in C^{n -\nu}$ hence the R.H.S is at least one time continously differentiable(the bottle neck being $a_{n -1}y^{(n -1)} \in C^{\min(k, 1)}$). Therefore $y^{(n)}\in C^1$. But this implies $y\in C^{n+1}$ hence we repeat the above argument to get $y^{(\nu)}\in C^{n+1 -\nu}$. The R.H.S has the bottleneck $a_{n -1}y^{(n -1)}\in C^{\min(k, n+1 -(n -1))}=C^{\min(k, 2)}$. We can repeat this procedure $m$ times and yield that $y^{(n)}\in C^{\min(k, m)}$. This means the best we can get is $y^{(n)}\in C^k$ or $y\in C^{n+k}$(after which the coeficient becomes the bottleneck).

\end{mystthmproof}
\begin{mystthmcorollary}{}{corollary-lineardifferentialoperator-5}If the coefficients are in $C^{\infty}$ then $y\in \mathrm{Null}(L)$ are in $C^{\infty}$.

\end{mystthmcorollary}
\subparagraph{Nullspace of $L$ is an $n$-dimensional subspace of $C^{\infty}$}

\paragraph{Linear Independence: Wronskian}

\paragraph{Solving $Ly=0$ with constant coefficients}

\paragraph{$Ly=0$ with general $a_\nu(x)$}

\begin{mystexercise}{1-dimensional Scattering Problem}{1dscat}\label{1dscat}Consider the 1d Schrodinger equation:
\begin{equation}
-\frac{d^2\psi}{dx^2}+V\psi = E\psi
\end{equation}
This is an eigenvalue problem of the operator $H= -\frac{d^2}{dx^2}+V$:
\begin{equation}
H\psi = E\psi
\end{equation}
Where $V$ has finite support $[ -a, a], a\in \mathbb{R}$. Denote:

\begin{enumerate}
\item $L:=\{x\in \mathbb{R}:x< -a\}$
\item $R:=\{x\in \mathbb{R}:x>a\}$
\end{enumerate}

\textbf{Problems}:

\begin{enumerate}
\item Show that in $L,R$, the general solution to the formal operator $H -E$ where $E=k^2$ is:

\begin{equation}
\begin{aligned}
\psi_k &=A\exp(ikx)+B\exp(-ikx) \quad x\in L  \\
\psi_k &=C\exp(ikx)+D\exp(-ikx) \quad x\in R
\end{aligned}
\end{equation}


\item Now consider two asymptotic boundary condition:

\begin{equation}
\begin{cases}
\Pi_L: D=0, x\to \infty \\
\Pi_R: C=0, x\to -\infty
\end{cases}
\end{equation}


where $\Pi_L$ depicts a left-incident wave with no incoming wave from the right and $\Pi_R$ is the opposite. Show that the solution(in $L,R$) for the BC $\Pi_1$ and $k>0$ is:

\begin{equation}
\begin{aligned}
\psi_k^L(x) = \begin{cases}
\exp(ikx)+ r_L(k)\exp(-ikx) & x\in L \\
t_L(k)\exp(ikx) & x\in R
\end{cases}
\end{aligned}
\end{equation}


Similarily, for $\Pi_2$, $k<0$:

\begin{equation}
\begin{aligned}
\psi_k^R(x)=\begin{cases}
t_R(k)\exp(ikx)& x\in L \\
\exp(ikx) + r_R(k)\exp(-ikx) & x\in R
\end{cases}
\end{aligned}
\end{equation}


where $t, r$ denotes the \textit{transmission} and \textit{reflection} coefficient.


\item Note that $\psi_k^*$ is also a solution. Use properties of Wronskian to shwo that:

\begin{equation}
|t_{L/R}|^2 + |r_{L/R}|^2=1
\end{equation}
\end{enumerate}

\end{mystexercise}
\begin{mystsolution}{Solution to Exercise~\ref{1dscat}}\begin{enumerate}
\item Fix a $k$. Consider the basis $\psi_k^L, {\psi_k^L}^*$. Form the Wronskian in $L$:

\begin{equation}
\begin{aligned}
W_L=\begin{vmatrix}
\exp(ikx)+r_L\exp(-ikx) &   \exp(-ikx)+r_L^*\exp(ikx) \\
ik\exp(ikx)-ikr_L\exp(-ikx) & -ik\exp(-ikx)+ikr_L^*\exp(ikx)
\end{vmatrix}
\end{aligned}
\end{equation}


Evaluate the determinant:

\begin{equation}
\begin{aligned}
W_L &= (-ik)\cancel{-ikr_L\exp(-2ikx) + ikr_L\exp(2ikx)}+ik|r_L|^2  \\
&-[ik\cancel{-ikr_K\exp(-2ikx)+ikr_L\exp(2ikx)}-ik|r_L|^2]  \\
&= 2ik(|r_L|^2-1)
\end{aligned}
\end{equation}


Next form Wronskian in $R$:

\begin{equation}
\begin{aligned}
W_R = \begin{vmatrix}
t_L\exp(ikx) & t_L\exp(-ikx) \\
ikt_L\exp(ikx) & -ikt_L\exp(-ikx) 
\end{vmatrix} = -2ik|t_L|^2
\end{aligned}
\end{equation}


Since we would expect $W_R=W_L=0$, henece:

\begin{equation}
\begin{aligned}
W_R-W_L&=0 \\
2ik(|r_L|^2+|t_L|^2-1) &=0
\end{aligned}
\end{equation}


Since $k>0$ we get as desired.
\end{enumerate}

4.Fix $k$. We repeat the same procedure with $\{\psi_k^L,\psi_{( -k)}^R\}$. We evaluate $W_L$:
\begin{equation}
\begin{aligned}
W_L &= \begin{vmatrix}
\exp(ikx)+r_L(k)\exp(-ikx) & t_R(-k)\exp(-ikx) \\
ik\exp(ikx)-ikr_L(k)\exp(-ikx) & -ikt_R(-k)\exp(-ikx)
\end{vmatrix} \\
&= -2ikt_R(-k)
\end{aligned}
\end{equation}
Similarily,
\begin{equation}
W_R = -2ikt_L(k)
\end{equation}
Again, $W_L=W_R=0$ so we get as desired.

\end{mystsolution}